% 第4章 標準難易度：基本4パターン・応用7パターンを各5題。
% 主分類・出題意図・別解の対応は docs/chapter4-standard-design-2026-10-04.md。

    \subsection{標準難易度}\label{節：標準難易度}\label{節：基本と応用から出題されるもの}
      ここでは,第3章の基本4パターンと応用7パターンから,それぞれ5題ずつ,計55題を出題する.
      同じ解法で解ける問題も,式の並び方や係数が変わると違って見える.
      まず移項や約分によって式を整理し,どの数量の差や比に注目すればよいかを考えよう.
      問題は型ごとにまとめずに並べてあるので,式の構造を見て解法を選んでほしい.
      方針と解法は問題と同じ問番号に対応する.
      別解のある問題では,最初の解法を理解したあとに,別の見方でも同じ答えになることを確認しよう.
      特に階比型と対数型には共通して使える見方がある.
      以下,\,$n$は正の整数とし,必要なときは
      \[
        S_n=\sum_{k=1}^{n}a_k
      \]
      とおく.

      \subsubsection{問題}\label{節：標準難易度：問題}\label{節：基本と応用から出題されるもの：問題}
        \begin{enumerate}[label=\textbf{問\arabic*.},leftmargin=3.8em,format=\bookitemlink{practice-basic}{problem}{answer}]
          % Q1
          \item $\displaystyle a_1=3$ とし,
                \[
                  (2a_n+3)a_{n+1}=3a_n
                \]
                を満たす数列$\{a_n\}$の一般項を求めよ.
          % S1
          \item $\displaystyle S_n=2a_n-n^2$
                                を満たす数列$\{a_n\}$の一般項を求めよ.
          % R1
          \item $\displaystyle a_1=3,\qquad a_{n+1}=2(n+1)a_n$
                で定められる数列の一般項を求めよ.
          % M1
          \item $\displaystyle a_1=2,\qquad a_{n+1}=2a_n+n-1$
                で定められる数列$\{a_n\}$の一般項を求めよ.
          % G1
          \item $\displaystyle a_1=6,\qquad a_{n+1}=-\frac12a_n$
                で定められる数列の一般項を求めよ.
          % T1
          \item $\displaystyle a_1=2,\,a_2=5$ とし,
                \[
                  a_{n+2}+6a_n=5a_{n+1}
                \]
                を満たす数列$\{a_n\}$の一般項を求めよ.
          % U1
          \item $\displaystyle a_1=3,\qquad b_1=1,$
                            \[
                              \begin{cases}
                                a_{n+1}=2a_n-b_n,\\
                                b_{n+1}=-a_n+2b_n
                              \end{cases}
                            \]
                                で定められる2つの数列$\{a_n\},\{b_n\}$の一般項を求めよ.
          % L1
          \item $\displaystyle a_1=4,\qquad a_{n+1}=\frac{a_n^3}{4}$
                で定められる数列$\{a_n\}$の一般項を求めよ.
          % C1
          \item $\displaystyle a_1=5,\qquad a_{n+1}=3a_n-4$
                で定められる数列$\{a_n\}$の一般項を求めよ.
          % D1
          \item $\displaystyle a_1=\frac56,\qquad 3a_{n+1}=3a_n-2$
                で定められる数列の一般項を求めよ.
          % F1
          \item $\displaystyle a_1=2,\qquad a_{n+1}=a_n+3n-1$
                で定められる数列の一般項を求めよ.
          % T2
          \item $\displaystyle a_1=3,\,a_2=1$ とし,
                \[
                  a_{n+2}+a_{n+1}=6a_n
                \]
                を満たす数列$\{a_n\}$の一般項を求めよ.
          % F2
          \item $\displaystyle a_1=1,\qquad a_{n+1}=a_n+3n^2+3n+1$
                で定められる数列の一般項を求めよ.
          % Q2
          \item $\displaystyle a_1=2,\,a_{n+1}=-\frac{2a_n}{3a_n+2}$
                で定められる数列$\{a_n\}$の一般項を求めよ.
          % C2
          \item $\displaystyle a_1=1,\qquad 2a_{n+1}+a_n=9$
                で定められる数列$\{a_n\}$の一般項を求めよ.
          % U2
          \item $\displaystyle a_1=3,\qquad b_1=1,$
                            \[
                              \begin{cases}
                                a_{n+1}=4a_n-2b_n,\\
                                b_{n+1}=a_n+b_n
                              \end{cases}
                            \]
                                で定められる2つの数列$\{a_n\},\{b_n\}$の一般項を求めよ.
          % D2
          \item $\displaystyle a_1=-4,\qquad 2a_{n+1}+a_n=3a_n+6$
                で定められる数列の一般項を求めよ.
          % G2
          \item $\displaystyle a_1=3,\qquad 2a_{n+1}-5a_n=a_{n+1}+a_n$
                で定められる数列の一般項を求めよ.
          % R2
          \item $\displaystyle a_1=6,\qquad
                 (n+1)(n+2)a_{n+1}=n(n+3)a_n$
                で定められる数列の一般項を求めよ.
          % M2
          \item $\displaystyle a_1=3,\qquad a_{n+1}+a_n=2n^2+2n+1$
                で定められる数列$\{a_n\}$の一般項を求めよ.
          % S2
          \item $\displaystyle S_n=a_n+(n-1)2^{n-1}$
                                を満たす数列$\{a_n\}$の一般項を求めよ.
          % L2
          \item $\displaystyle a_1=5,\qquad a_{n+1}=a_n^2-4a_n+6$
                で定められる数列$\{a_n\}$の一般項を求めよ.
          % T3
          \item $\displaystyle a_1=1,\,a_2=0$ とし,
                \[
                  a_{n+2}-4a_{n+1}+4a_n=0
                \]
                を満たす数列$\{a_n\}$の一般項を求めよ.
          % M3
          \item $\displaystyle a_1=1,\qquad a_{n+1}=3a_n-2^n$
                で定められる数列$\{a_n\}$の一般項を求めよ.
          % U3
          \item $\displaystyle a_1=1,\qquad b_1=2,$
                            \[
                              \begin{cases}
                                a_{n+1}=2a_n+b_n,\\
                                b_{n+1}=2b_n
                              \end{cases}
                            \]
                                で定められる2つの数列$\{a_n\},\{b_n\}$の一般項を求めよ.
          % Q3
          \item $\displaystyle a_1=4,\,a_{n+1}=\frac{3a_n+2}{a_n+2}$
                で定められる数列$\{a_n\}$の一般項を求めよ.
          % G3
          \item $\displaystyle a_1=4,\qquad (n+1)a_{n+1}=\frac{n+1}{3}a_n$
                で定められる数列の一般項を求めよ.
          % D3
          \item $\displaystyle a_1=1-\sqrt2,\qquad a_{n+1}+n=a_n+n+\sqrt2$
                で定められる数列の一般項を求めよ.
          % L3
          \item $\displaystyle a_1=2,\qquad a_{n+1}=a_n^{\frac{n+2}{n}}$
                で定められる正の数列$\{a_n\}$の一般項を求めよ.
          % S3
          \item $\displaystyle S_1=1,\qquad S_{n+1}=3S_n+2$
                                を満たす数列$\{a_n\}$の一般項を求めよ.
          % F3
          \item $\displaystyle a_1=3,\qquad a_{n+1}=a_n+2\cdot3^n$
                で定められる数列の一般項を求めよ.
          % R3
          \item $\displaystyle a_1=4,\qquad (n+1)^2a_{n+1}=n^2a_n$
                で定められる数列の一般項を求めよ.
          % C3
          \item $\displaystyle a_1=0,\qquad a_{n+1}-a_n=\frac{6-a_n}{3}$
                で定められる数列$\{a_n\}$の一般項を求めよ.
          % S4
          \item $\displaystyle a_1=1,\qquad a_{n+1}=2S_n+2^n$
                                で定められる数列$\{a_n\}$の一般項を求めよ.
          % Q4
          \item $\displaystyle a_1=1,\,a_{n+1}=2+\frac3{a_n}$
                で定められる数列$\{a_n\}$の一般項を求めよ.
          % G4
          \item 数列$\{a_n\}$は,ある実数$r$について
                \[
                  a_{n+1}=ra_n
                \]
                を満たし,\,$a_2=-6$,\, $a_5=48$である.\,$r$および一般項$a_n$を求めよ.
          % F4
          \item $\displaystyle a_1=0,\qquad a_{n+1}=a_n+\frac{2}{(n+1)(n+3)}$
                で定められる数列の一般項を求めよ.
          % T4
          \item $\displaystyle a_1=3,\,a_2=5$ とし,
                \[
                  a_{n+2}=2a_{n+1}+3a_n-8
                \]
                を満たす数列$\{a_n\}$の一般項を求めよ.
          % U4
          \item $\displaystyle a_1=0,\qquad b_1=1,$
                            \[
                              \begin{cases}
                                a_{n+1}=2a_n+b_n+1,\\
                                b_{n+1}=a_n+2b_n-1
                              \end{cases}
                            \]
                                で定められる2つの数列$\{a_n\},\{b_n\}$の一般項を求めよ.
          % L4
          \item $\displaystyle a_1=1,\qquad a_{n+1}=2\sqrt{a_n}$
                で定められる数列$\{a_n\}$の一般項を求めよ.
          % C4
          \item 実数$t$に対し,
                \[
                  a_1=t,\qquad a_{n+1}=4a_n-9
                \]
                で数列$\{a_n\}$を定める.一般項を求め,この数列が定数数列となる$t$の値を求めよ.
          % R4
          \item $\displaystyle a_1=2,\qquad a_{n+1}=2^{2n+1}a_n$
                で定められる数列の一般項を求めよ.
          % D4
          \item 数列$\{a_n\}$は,ある実数$d$について
                \[
                  a_{n+1}=a_n+d
                \]
                を満たし,\,$a_3=8$,\, $a_8=-7$である.\,$d$,\, $a_1$および一般項$a_n$を求めよ.
          % M4
          \item $\displaystyle a_1=2,\qquad a_{n+1}=2a_n+3\cdot2^n$
                で定められる数列$\{a_n\}$の一般項を求めよ.
          % U5
          \item $\displaystyle a_1=1,\qquad b_1=0,$
                            \[
                              \begin{cases}
                                a_{n+1}=a_n+2b_n,\\
                                b_{n+1}=a_n
                              \end{cases}
                            \]
                                で定められる2つの数列$\{a_n\},\{b_n\}$の一般項を求めよ.
          % R5
          \item $\displaystyle a_1=1,\qquad a_{n+1}=\frac{n(n+1)}2a_n$
                で定められる数列の一般項を求めよ.
          % Q5
          \item $\displaystyle a_1=2,\,a_{n+1}=2-\frac1{a_n}$
                で定められる数列$\{a_n\}$の一般項を求めよ.
          % G5
          \item $\displaystyle a_1=0,\qquad a_{n+1}=7a_n$
                で定められる数列の一般項を求めよ.
          % L5
          \item $\displaystyle a_1=4,\qquad a_{n+1}=\frac{a_n^2}{2^{n-1}}$
                で定められる数列$\{a_n\}$の一般項を求めよ.
          % F5
          \item $\displaystyle a_1=1,\qquad a_{n+1}=a_n+(-1)^n$
                で定められる数列の一般項を求めよ.
          % C5
          \item 実数$p,q$を用いて,
                \[
                  a_{n+1}=pa_n+q
                \]
                で定められる数列$\{a_n\}$について,
                \[
                  a_1=8,\qquad a_2=5,\qquad a_3=\frac72
                \]
                であることが分かっている.\,$p,q$の値と一般項を求めよ.
          % T5
          \item $\displaystyle a_1=1,\,a_2=2$ とし,
                \[
                  a_{n+2}-3a_{n+1}+2a_n=1
                \]
                を満たす数列$\{a_n\}$の一般項を求めよ.
          % D5
          \item $\displaystyle a_1=\frac74,\qquad 4a_{n+1}-4a_n=0$
                で定められる数列の一般項を求めよ.
          % M5
          \item $\displaystyle a_1=0,\qquad a_{n+1}=2a_n+3(-1)^n$
                で定められる数列$\{a_n\}$の一般項を求めよ.
          % S5
          \item $\displaystyle a_1=2,\qquad S_n=\frac{n+2}{3}a_n$
                                を満たす数列$\{a_n\}$の一般項を求めよ.
        \end{enumerate}

      \subsubsection{方針}\label{節：標準難易度：方針}\label{節：基本と応用から出題されるもの：方針}
        \begin{enumerate}[label=\textbf{問\arabic*.},leftmargin=3.8em,format=\bookitemlink{practice-basic}{hint}{problem}]
          % Q1
          \item $a_{n+1}$について解き,両辺の逆数を取る.\,$\displaystyle \frac{1}{a_n}$の増え方に注目する.
          % S1
          \item $n=1$から初項を求める. 次に$S_{n+1}-S_n=a_{n+1}$を使い, 総和を消す. 得られた漸化式では$a_n$に$n$の一次式を加えて等比数列を作る.
          % R1
          \item $(n+1)!=(n+1)n!$を使い,\,$a_n$を$n!$で割る.
          % M1
          \item $a_n$から適切な一次式を引き,残りが等比数列になるようにする.
          % G1
          \item 初項から第$n$項までに,$-\dfrac12$を掛ける回数を考える.
          % T1
          \item 特性方程式を解き,\,$a_{n+1}-2a_n$と$a_{n+1}-3a_n$を求める.二式の差で$a_n$が残る.
          % U1
          \item 2つの式を足すと和は変わらず, 引くと差は3倍になる. 和と差を求めてから$a_n,b_n$に戻す.
          % L1
          \item 全項が正であることを確かめ,底を$2$とする対数を取る.
          % C1
          \item $a_n$の値が変わらないとしたときの値を求め,その値との差を考える.
          % D1
          \item 両辺を$3$で割り,\,$a_{n+1}-a_n$を求める.
          % F1
          \item $a_{k+1}-a_k=3k-1$を$k=1$から$n-1$まで加える.
          % T2
          \item 特性方程式の解は$2$と$-3$である.負の解に対応する補助数列と$(-3)^{n-1}$の符号に注意する.
          % F2
          \item $3n^2+3n+1$が$(n+1)^3-n^3$であることを使う.
          % Q2
          \item 逆数を取って一次の漸化式に帰着させる.最初の数項を計算すると,より短い解法も見つかる.
          % C2
          \item まず$a_{n+1}$について解く.不動点との差は負の公比をもつ等比数列になる.
          % U2
          \item $a_n+t b_n$が等比数列になるように$t$を選ぶ. $t=-1,-2$で2本の等比数列を得られる. 片方を消して隣接3項間の漸化式にする別解もある.
          % D2
          \item 両辺から$a_n$を引き,その後で$2$で割る.
          % G2
          \item $a_{n+1}$を左辺に,\,$a_n$を右辺にまとめる.
          % R2
          \item 倍率を$\displaystyle \frac{n}{n+1}\frac{n+3}{n+2}$と分解し,順に掛ける.
          % M2
          \item 右辺を$(n+1)^2+n^2$と見れば,引くべき二次式が分かる.
          % S2
          \item $S_{n+1}-S_n=a_{n+1}$に代入する. 両辺の$a_{n+1}$が消え,\,$a_n$が直接求まる. $n=1$の与式だけでは初項が決まらないことにも注意する.
          % L2
          \item 平方完成すると$a_{n+1}-2=(a_n-2)^2$となる.正の量$a_n-2$の対数を取る.
          % T3
          \item 特性方程式は重解$2$を持つ.\,$a_{n+1}-2a_n$を等比数列にした後,\,$\displaystyle \frac{a_n}{2^{n-1}}$を考える.
          % M3
          \item $2^n$に関係する等比数列を引くか,$2^{n-1}$で割って一次の漸化式にする.
          % U3
          \item $b_n$は$a_n$と関係なく先に解ける. その結果を第1式に代入し,\,$\displaystyle \frac{a_n}{2^{n-1}}$の階差を調べる.
          % Q3
          \item 不動点は$2$と$-1$である.\,$\displaystyle \frac{a_n-2}{a_n+1}$を考える.\,$\displaystyle \frac{1}{a_n-2}$を考える別解もある.
          % G3
          \item $n$は正整数なので,\,$n+1$で割ることができる.
          % D3
          \item 両辺に現れる$n$を消去してから,隣り合う項の差を見る.
          % L3
          \item 対数を取った数列の階比が$\displaystyle \frac{n+2}{n}$となる.積の途中の因子が消えることに注目する.
          % S3
          \item $S_n+1$を等比数列として解いてから,\,$a_n=S_n-S_{n-1}$を使う. この差の公式を使う範囲と$a_1$を区別する.
          % F3
          \item 階差を加えた後,\,$3+3^2+\cdots+3^{n-1}$を求める.
          % R3
          \item 左辺は第$n+1$項に$(n+1)^2$を掛けた形,右辺は第$n$項に$n^2$を掛けた形である.
          % C3
          \item 右辺が$0$になる値$6$との差を取ると,等比型になる.
          % S4
          \item $S_{n+1}=S_n+a_{n+1}$から$S_n$の漸化式を作る. $S_n+2^n$が等比数列になる. 1つ前の式との差から$a_n$の漸化式を作る別解もある.
          % Q4
          \item 右辺を一つの分数にまとめる.不動点$3,-1$からの比を取ると,公比が負の等比数列になる.
          % G4
          \item 第$2$項から第$5$項までには$r$を$3$回掛ける.
          % F4
          \item $\displaystyle \frac{2}{(k+1)(k+3)}=\frac1{k+1}-\frac1{k+3}$と分解する.
          % T4
          \item 定数解$2$を引く.その後の特性方程式の一つの解$3$を使うと,初項が$0$の補助数列が現れる.
          % U4
          \item 定数項にも注意して2つの式を足し引きする. 和は3倍され, 差には2が加えられる.
          % L4
          \item 平方根を$\displaystyle a_n^{\frac{1}{2}}$と書き,底を$2$とする対数を取る.
          % C4
          \item 不動点との差を取る.その差の初項が$0$となる場合も,分けて確認する.
          % R4
          \item $2$の累乗を掛けると指数が加わる.対数を取って和に直してもよい.
          % D4
          \item 第$3$項から第$8$項までには,公差を何回加えるかを数える.
          % M4
          \item $2^{n-1}$で割った数列を考える.変形後は等差型になる.
          % U5
          \item 第1式の添字を1つ進め, 第2式を代入して$b_n$を消す. $a_2$はもとの連立から求める. 最後に$b_n$を求めるときは初項も確認する.
          % R5
          \item $(n-1)!n!$の次の項との比が$n(n+1)$であることを使う.
          % Q5
          \item 不動点の方程式は重解$1$を持つ.二つの差の比の代わりに,\,$\displaystyle \frac{1}{a_n-1}$を考える.
          % G5
          \item $0$に$7$を掛けても$0$のままである.項で割らずに考える.
          % L5
          \item 対数を取ると$b_{n+1}=2b_n-n+1$となる.この式を満たす一次式を探す.
          % F5
          \item $(-1)^k$の和は等比数列の和として求められる.最初の数項も調べる.
          % C5
          \item 最初の2回の更新式を引けば$q$が消える.係数を求めた後,不動点との差を取る.
          % T5
          \item 定数を引くだけでは右辺を$0$にできない.\,$d_n=a_{n+1}-a_n$を考え,一次の漸化式に帰着させる.
          % D5
          \item $a_{n+1}=a_n$なら,項の値は変わらない.
          % M5
          \item $(-1)^n$も指数の形である.$(-1)^{n-1}$で割った数列を考える.
          % S5
          \item $S_{n+1}-S_n=a_{n+1}$から$\displaystyle \frac{a_{n+1}}{a_n}$を求める. 比を順に掛けると, 分子と分母の大部分が打ち消し合う.
        \end{enumerate}

      \subsubsection{解答}\label{節：標準難易度：解答}\label{節：基本と応用から出題されるもの：解答}
        \begin{enumerate}[label=\textbf{問\arabic*.},leftmargin=3.8em,format=\bookitemlink{practice-basic}{answer}{problem}]
          % Q1
          \item \[
                  a_n=\frac{3}{2n-1}
                \]
          % S1
          \item \[
                  a_n=6\cdot2^{n-1}-2n-3
                \]
          % R1
          \item \[
                  a_n=3\cdot2^{n-1}n!
                \]
          % M1
          \item \[
                  a_n=3\cdot2^{n-1}-n
                \]
          % G1
          \item \[
                  a_n=6\left(-\frac12\right)^{n-1}
                \]
          % T1
          \item \[
                  a_n=2^{n-1}+3^{n-1}
                \]
          % U1
          \item \[
                  a_n=2+3^{n-1},\qquad b_n=2-3^{n-1}
                \]
          % L1
          \item \[
                  a_n=2^{\,1+3^{n-1}}
                \]
          % C1
          \item \[
                  a_n=2+3^n
                \]
          % D1
          \item \[
                  a_n=\frac56-\frac23(n-1)=\frac{9-4n}{6}
                \]
          % F1
          \item \[
                  a_n=\frac{3n^2-5n+6}{2}
                \]
          % T2
          \item \[
                  a_n=2^n+(-3)^{n-1}
                \]
          % F2
          \item \[
                  a_n=n^3
                \]
          % Q2
          \item \[
                  a_n=\frac4{5(-1)^{n-1}-3}
                \]
          % C2
          \item \[
                  a_n=3-2\left(-\frac12\right)^{n-1}
                \]
          % U2
          \item \[
                  a_n=4\cdot3^{n-1}-2^{n-1},\qquad b_n=2\cdot3^{n-1}-2^{n-1}
                \]
          % D2
          \item \[
                  a_n=3n-7
                \]
          % G2
          \item \[
                  a_n=3\cdot6^{n-1}
                \]
          % R2
          \item \[
                  a_n=\frac{2(n+2)}n
                \]
          % M2
          \item \[
                  a_n=n^2+2(-1)^{n-1}
                \]
          % S2
          \item \[
                  a_n=(n+1)2^{n-1}
                \]
          % L2
          \item \[
                  a_n=2+3^{\,2^{n-1}}
                \]
          % T3
          \item \[
                  a_n=(2-n)2^{n-1}
                \]
          % M3
          \item \[
                  a_n=2^n-3^{n-1}
                \]
          % U3
          \item \[
                  a_n=n2^{n-1},\qquad b_n=2^n
                \]
          % Q3
          \item \[
                  a_n=\frac{10\cdot4^{n-1}+2}{5\cdot4^{n-1}-2}
                \]
          % G3
          \item \[
                  a_n=4\left(\frac13\right)^{n-1}
                \]
          % D3
          \item \[
                  a_n=1+(n-2)\sqrt2
                \]
          % L3
          \item \[
                  a_n=2^{\frac{n(n+1)}2}
                \]
          % S3
          \item \[
                  a_1=1,\qquad a_n=4\cdot3^{n-2}\quad(n\ge2)
                \]
          % F3
          \item \[
                  a_n=3^n
                \]
          % R3
          \item \[
                  a_n=\frac4{n^2}
                \]
          % C3
          \item \[
                  a_n=6\left\{1-\left(\frac23\right)^{n-1}\right\}
                \]
          % S4
          \item \[
                  a_n=2\cdot3^{n-1}-2^{n-1}
                \]
          % Q4
          \item \[
                  a_n=\frac{3^n+(-1)^n}{3^{n-1}-(-1)^n}
                \]
          % G4
          \item \[
                  r=-2,\qquad a_n=3(-2)^{n-1}
                \]
          % F4
          \item \[
                  a_n=\frac56-\frac1{n+1}-\frac1{n+2}
                \]
          % T4
          \item \[
                  a_n=2+3^{n-1}
                \]
          % U4
          \item \[
                  a_n=\frac{3^{n-1}+2n-3}{2},\qquad b_n=\frac{3^{n-1}-2n+3}{2}
                \]
          % L4
          \item \[
                  a_n=2^{\,2-2^{2-n}}
                \]
          % C4
          \item \[
                  a_n=3+(t-3)4^{n-1},\qquad
                \text{定数数列となるのは }t=3
                \]
          % R4
          \item \[
                  a_n=2^{n^2}
                \]
          % D4
          \item \[
                  d=-3,\qquad a_1=14,\qquad a_n=17-3n
                \]
          % M4
          \item \[
                  a_n=(3n-1)2^{n-1}
                \]
          % U5
          \item \[
                  a_n=\frac{2^n+(-1)^{n-1}}{3},\qquad b_n=\frac{2^{n-1}-(-1)^{n-1}}{3}
                \]
          % R5
          \item \[
                  a_n=\frac{(n-1)!\,n!}{2^{n-1}}
                \]
          % Q5
          \item \[
                  a_n=\frac{n+1}{n}
                \]
          % G5
          \item \[
                  a_n=0
                \]
          % L5
          \item \[
                  a_n=2^{\,n+2^{n-1}}
                \]
          % F5
          \item \[
                  a_n=\frac{1+(-1)^{n-1}}2
                \]
          % C5
          \item \[
                  p=\frac12,\qquad q=1,\qquad
                a_n=2+6\left(\frac12\right)^{n-1}
                \]
          % T5
          \item \[
                  a_n=2^n-n
                \]
          % D5
          \item \[
                  a_n=\frac74
                \]
          % M5
          \item \[
                  a_n=(-1)^{n-1}-2^{n-1}
                \]
          % S5
          \item \[
                  a_n=n(n+1)
                \]
        \end{enumerate}

      \subsubsection{解法}\label{節：標準難易度：解法}\label{節：基本と応用から出題されるもの：解法}
        \begin{enumerate}[label=\textbf{問\arabic*.},leftmargin=3.8em,format=\bookitemlink{practice-basic}{solution}{problem}]
          % Q1
          \item $a_1>0$であり,\,$a_n>0$なら
                \[
                  a_{n+1}=\frac{3a_n}{2a_n+3}>0
                \]
                である.したがって,すべての項は正で,分母$2a_n+3$も$0$ではない.
                両辺の逆数を取ると,
                \[
                  \frac1{a_{n+1}}=\frac{2a_n+3}{3a_n}
                  =\frac1{a_n}+\frac23.
                \]
                $\displaystyle b_n=\frac{1}{a_n}$と置けば,初項$\displaystyle \frac{1}{3}$,公差$\displaystyle \frac{2}{3}$の等差数列になる.よって,
                \[
                  b_n=\frac13+\frac23(n-1)=\frac{2n-1}{3},
                  \qquad
                  a_n=\frac3{2n-1}.
                \]
                求めた式もすべての正整数$n$で正であり,逆数を戻す際の分母も$0$にならない.

                \noindent \textbf{【別解】}\\
                最初の数項は
                \[
                  a_1=\frac31,\quad a_2=\frac33,\quad a_3=\frac35
                \]
                なので,\,$\displaystyle a_n=\frac{3}{2n-1}$と予想できる.
                この式は$n=1$で成り立つ.\,$\displaystyle a_k=\frac{3}{2k-1}$と仮定すると,
                \[
                  a_{k+1}
                  =\frac{3\cdot\frac3{2k-1}}{2\cdot\frac3{2k-1}+3}
                  =\frac3{2k+1}
                  =\frac3{2(k+1)-1}.
                \]
                したがって,数学的帰納法によって一般項が確かめられる.
          % S1
          \item $n=1$を代入すると$a_1=2a_1-1$であるから,\,$a_1=1$である.
                                また,
                            \begin{align*}
                              a_{n+1}
                                &=S_{n+1}-S_n\\
                                &=2a_{n+1}-(n+1)^2\\
                                &\quad{}-2a_n+n^2
                            \end{align*}
                                より,
                            \[
                              a_{n+1}=2a_n+2n+1.
                            \]
                                右辺に$n$があるので, 定数を足すだけでは等比数列にならない.
                                $b_n=a_n+2n+3$とおくと,
                            \begin{align*}
                              b_{n+1}
                                &=2a_n+2n+1+2(n+1)+3\\
                                &=2(a_n+2n+3)=2b_n.
                            \end{align*}
                                $b_1=1+2+3=6$より$b_n=6\cdot2^{n-1}$である.
                                したがって,
                            \[
                              a_n=6\cdot2^{n-1}-2n-3.
                            \]
          % R1
          \item 係数$n+1$は階乗に含まれるので,漸化式の両辺を$(n+1)!$で割ると,
                \[
                  \frac{a_{n+1}}{(n+1)!}=2\frac{a_n}{n!}
                \]
                となる.そこで
                \[
                  b_n=\frac{a_n}{n!}
                \]
                と置くと,\,$b_{n+1}=2b_n$,\, $b_1=3$である.
                よって,\,$b_n=3\cdot2^{n-1}$であり,
                \[
                  a_n=3\cdot2^{n-1}n!
                \]
                を得る.分母を$(n-1)!$にするのではなく,係数$n+1$と添字が合う$n!$を選ぶことがポイントである.

                \noindent \textbf{【別解】}\\
                倍率を初項から順に掛けると,
                \[
                  a_n
                    =3\prod_{k=1}^{n-1}2(k+1)
                    =3\cdot2^{n-1}(2\cdot3\cdots n)
                    =3\cdot2^{n-1}n!
                \]
                である.この積には$2$が$n-1$個含まれる.
                $n=1$では空積を$1$とするので,同じ式が成り立つ.
          % M1
          \item 定数を引くだけでは$n-1$を消せないので,同じ漸化式を満たす一次式$g(n)=un+v$を探す.
                必要な条件は,
                \[
                  g(n+1)=2g(n)+n-1
                \]
                である.これに$g(n)=un+v$を代入して係数を比べると,
                \[
                  u=2u+1,\qquad u+v=2v-1
                \]
                なので,\,$u=-1,v=0$を得る.つまり$g(n)=-n$である.
                \,$b_n=a_n-g(n)=a_n+n$とおくと,
                \[
                  b_{n+1}=a_{n+1}+n+1=2a_n+2n=2b_n
                \]
                であり,\,$b_1=2+1=3$だから,
                \[
                  b_n=3\cdot2^{n-1},\qquad
                  a_n=3\cdot2^{n-1}-n
                \]
                となる.最終式に$n=1$を代入すると$2$である.

                \noindent \textbf{【別解：階差を先に求める】}\\
                隣り合う更新式を引くと,
                \[
                  d_n=a_{n+1}-a_n,\qquad d_{n+1}=2d_n+1
                \]
                となる.\,$a_2=4$より$d_1=2$なので,
                \[
                  d_n+1=3\cdot2^{n-1}
                \]
                である.従って,\,$n\ge2$では,
                \[
                  a_n=2+\sum_{k=1}^{n-1}(3\cdot2^{k-1}-1)
                     =3\cdot2^{n-1}-n
                \]
                となる.\,$n=1$でも最終式は成り立つ.
          % G1
          \item 次の項は前の項に$-\dfrac12$を掛けて得られる.
                したがって,初項$6$,\,公比$-\dfrac12$の等比数列であるから,
                \[
                  a_n=6\left(-\frac12\right)^{n-1}
                \]
                である.最初の数項は$6,\,-3,\,\dfrac32,\,-\dfrac34,\ldots$となる.
                公比が負であるため符号は交互に変わり,その絶対値は毎回$\dfrac12$倍になる.
          % T1
          \item 漸化式を$a_{n+2}-5a_{n+1}+6a_n=0$と整理する.
                特性方程式は
                \[
                  x^2-5x+6=(x-2)(x-3)=0
                \]
                なので,二つの解は$2$と$3$である.
                そこで,
                \[
                  b_n=a_{n+1}-2a_n,\qquad c_n=a_{n+1}-3a_n
                \]
                と置くと,
                \begin{align*}
                  b_{n+1}
                  &=a_{n+2}-2a_{n+1}
                  \\
                  &=3(a_{n+1}-2a_n)=3b_n,\\
                  c_{n+1}
                  &=a_{n+2}-3a_{n+1}
                  \\
                  &=2(a_{n+1}-3a_n)=2c_n.
                \end{align*}
                初項は$b_1=5-2\cdot2=1,\,c_1=5-3\cdot2=-1$だから,
                \[
                  b_n=3^{n-1},\qquad c_n=-2^{n-1}.
                \]
                二式の差を取ると$b_n-c_n=a_n$である.よって,
                \[
                  a_n=2^{n-1}+3^{n-1}.
                \]
                二つの等比数列から,共通して含まれている$a_{n+1}$を消去することが要点である.
          % U1
          \item $s_n=a_n+b_n,\ d_n=a_n-b_n$とおく.
                                2つの式を足し引きすると,
                            \[
                              s_{n+1}=s_n,\qquad d_{n+1}=3d_n.
                            \]
                                $s_1=3+1=4,\ d_1=3-1=2$より,
                            \[
                              s_n=4,\qquad d_n=2\cdot3^{n-1}.
                            \]
                                和と差からもとの2数列を復元すると,
                            \begin{align*}
                              a_n&=\frac{s_n+d_n}{2}=2+3^{n-1},\\
                              b_n&=\frac{s_n-d_n}{2}=2-3^{n-1}.
                            \end{align*}

                \noindent \textbf{【別解】}\\
                和が変わらないことから$a_n+b_n=4$である.
                                したがって$b_n=4-a_n$を第1式に代入して,
                            \[
                              a_{n+1}=2a_n-(4-a_n)=3a_n-4.
                            \]
                                $a_{n+1}-2=3(a_n-2)$と$a_1-2=1$より$a_n-2=3^{n-1}$である.
                                $b_n=4-a_n$を使えば, 2数列を求められる.
          % L1
          \item $a_1=4>0$であり,\,$a_n>0$なら$\displaystyle a_{n+1}=\frac{a_n^3}{4}>0$である.
                従って数学的帰納法により全ての項が正となり,対数を取ることができる.
                \,$b_n=\log_2a_n$とおくと,
                \[
                  b_{n+1}=3b_n-2,\qquad b_1=2
                \]
                となる.不動点$1$との差を取れば,
                \[
                  b_{n+1}-1=3(b_n-1)
                \]
                であり,\,$b_1-1=1$だから,
                \[
                  b_n=1+3^{n-1}
                \]
                を得る.\,$b_n=\log_2a_n$をもとの数列に戻して,
                \[
                  a_n=2^{\,1+3^{n-1}}
                \]
                となる.

                \noindent \textbf{【別解：指数を直接追う】}\\
                \,$\displaystyle c_n=\frac{a_n}{2}$とおくと,
                \[
                  c_{n+1}=\frac{a_n^3}{8}=c_n^3,\qquad c_1=2
                \]
                となる.1回進むたびに指数が$3$倍されるので,
                \[
                  c_n=2^{3^{n-1}}
                \]
                である.実際,\,$n=1$では$2$となり,この式の両辺を3乗すれば$n+1$の式が得られる.
                従って,\,$a_n=2c_n=2^{1+3^{n-1}}$となる.

                \noindent \textbf{【別解：階比を補助数列にする】}\\
                全項が正なので,\,$\displaystyle r_n=\frac{a_{n+1}}{a_n}=\frac{a_n^2}{4}$とおける.
                この階比について,
                \[
                  r_{n+1}=\frac{a_{n+1}^2}{4}
                         =\frac{a_n^6}{64}=r_n^3,\qquad r_1=4
                \]
                が成り立つ.指数を追うと$r_n=4^{3^{n-1}}$となるから,\,$n\ge2$では,
                \[
                  a_n=4\prod_{k=1}^{n-1}r_k
                     =4^{\,1+\sum_{k=1}^{n-1}3^{k-1}}
                     =4^{\,\frac{1+3^{n-1}}{2}}
                     =2^{\,1+3^{n-1}}
                \]
                である.ここでは階比の積を,指数の和に読み替えている.
                最終式は$n=1$でも成り立つ.
          % C1
          \item 定数項$-4$を消すため,同じ漸化式を満たす定数$\alpha$を探す.
                \[
                  \alpha=3\alpha-4\quad\Longleftrightarrow\quad\alpha=2
                \]
                この式をもとの漸化式から引くと,
                \[
                  a_{n+1}-2=3(a_n-2)
                \]
                となる.そこで,\,$b_n=a_n-2$とおくと,
                \[
                  b_{n+1}=3b_n,\qquad b_1=5-2=3
                \]
                である.初項が$3$,公比が$3$の等比数列なので,
                \[
                  b_n=3\cdot3^{n-1}=3^n
                \]
                従って,
                \[
                  a_n=2+3^n
                \]
                となる.\,$n=1$を代入すると$5$となり,初期条件も満たす.

                \noindent \textbf{【別解：階差に注目する】}\\
                漸化式を1つずらして引くと,
                \[
                  a_{n+2}-a_{n+1}=3(a_{n+1}-a_n)
                \]
                である.\,$d_n=a_{n+1}-a_n$とおくと,\,$a_2=11$より$d_1=6$であり,
                \[
                  d_n=6\cdot3^{n-1}
                \]
                となる.従って,\,$n\ge2$では,
                \[
                  a_n=5+6\sum_{k=1}^{n-1}3^{k-1}
                     =5+3(3^{n-1}-1)=2+3^n
                \]
                を得る.得られた式は$n=1$でも成り立つ.
          % D1
          \item 両辺を$3$で割ると,
                \[
                  a_{n+1}=a_n-\frac23
                \]
                となる.隣り合う項の差が常に$-\dfrac23$であるから,これは公差$-\dfrac23$の等差数列である.
                初項から第$n$項までにこの公差を$n-1$回加えるので,
                \[
                  a_n=\frac56-\frac23(n-1)=\frac{9-4n}{6}
                \]
                である.この式に$n=1$を代入すると,\,$a_1=\dfrac56$も確かめられる.
          % F1
          \item 隣り合う項の差は$3n-1$なので,初項にこれまでの増加分を加えればよい.
                $n\ge2$のとき,
                \begin{align*}
                  a_n
                    &=2+\sum_{k=1}^{n-1}(3k-1)\\
                    &=2+3\frac{n(n-1)}2-(n-1)\\
                    &=\frac{3n^2-5n+6}{2}
                \end{align*}
                となる.最後の式は$n=1$でも$2$となるから,すべての正整数$n$について成り立つ.
                和の上限を$n$にしてしまうと,第$n+1$項を求めることになるので注意する.
          % T2
          \item 特性方程式は
                \[
                  x^2+x-6=(x-2)(x+3)=0
                \]
                なので,二つの解は$2$と$-3$である.
                \[
                  b_n=a_{n+1}-2a_n,\qquad c_n=a_{n+1}+3a_n
                \]
                と置くと,
                \begin{align*}
                  b_{n+1}
                  &=a_{n+2}-2a_{n+1}
                  \\
                  &=-3(a_{n+1}-2a_n)=-3b_n,\\
                  c_{n+1}
                  &=a_{n+2}+3a_{n+1}
                  \\
                  &=2(a_{n+1}+3a_n)=2c_n.
                \end{align*}
                $b_1=1-2\cdot3=-5,\,c_1=1+3\cdot3=10$だから,
                \[
                  b_n=-5(-3)^{n-1},\qquad c_n=10\cdot2^{n-1}.
                \]
                $c_n-b_n=5a_n$なので,
                \[
                  5a_n=10\cdot2^{n-1}+5(-3)^{n-1}.
                \]
                よって,
                \[
                  a_n=2^n+(-3)^{n-1}.
                \]
                負の公比も等比数列の公比としてそのまま扱える.
          % F2
          \item 二項定理より,
                \[
                  3n^2+3n+1=(n+1)^3-n^3
                \]
                である.したがって,\,$n\ge2$のとき,
                \begin{align*}
                  a_n
                    &=1+\sum_{k=1}^{n-1}\{(k+1)^3-k^3\}\\
                    &=1+(2^3-1^3)+(3^3-2^3)\\
                    &\quad{}+\cdots+\{n^3-(n-1)^3\}\\
                    &=1+n^3-1\\
                    &=n^3
                \end{align*}
                となる.途中の立方数が次々に相殺するため,平方和を別に計算する必要はない.
                $n=1$でも$a_1=1=1^3$なので,結論はすべての正整数$n$で成り立つ.

                \noindent \textbf{【別解1】}\\
                最初の数項を計算すると$1,\,8,\,27,\ldots$となり,\,$a_n=n^3$が予想できる.
                これを数学的帰納法で確かめる.
                $n=1$では$a_1=1=1^3$である.
                $a_n=n^3$を仮定すると,
                \[
                  a_{n+1}=n^3+3n^2+3n+1=(n+1)^3
                \]
                となるから,予想はすべての正整数$n$について正しい.

                \noindent \textbf{【別解2】}\\
                漸化式を
                \[
                  a_{n+1}-(n+1)^3=a_n-n^3
                \]
                と書けば,\,$a_n-n^3$は一定である.
                その値は$a_1-1^3=0$だから,\,$a_n=n^3$となる.
          % Q2
          \item 初項から計算すると,
                \[
                  a_1=2,\qquad a_2=-\frac12,\qquad a_3=2.
                \]
                $2$と$\displaystyle -\frac{1}{2}$は互いに移り合い,これらの値で分母$3a_n+2$はそれぞれ$8$と$\displaystyle \frac{1}{2}$になる.
                したがって,数列はすべての正整数$n$について定まり,どの項も$0$ではない.
                一般項を一つの式で表すために,\,$\displaystyle b_n=\frac{1}{a_n}$と置くと,
                \[
                  b_{n+1}=-b_n-\frac32,\qquad b_1=\frac12.
                \]
                定数項を消すために$\displaystyle \frac{3}{4}$を加えれば,
                \[
                  b_{n+1}+\frac34=-\left(b_n+\frac34\right).
                \]
                よって,
                \[
                  b_n+\frac34=\left(\frac12+\frac34\right)(-1)^{n-1}
                  =\frac54(-1)^{n-1}.
                \]
                したがって,
                \[
                  b_n=\frac{5(-1)^{n-1}-3}{4},
                  \qquad
                  a_n=\frac4{5(-1)^{n-1}-3}.
                \]
                この分母は$n$が奇数なら$2$,偶数なら$-8$なので,\,$0$にはならない.

                \noindent \textbf{【別解】}\\
                漸化式の右辺を$\displaystyle f(x)=\frac{-2x}{3x+2}$と書くと,
                \[
                  f(2)=-\frac12,\qquad f\left(-\frac12\right)=2.
                \]
                したがって,初項からこの二つの値を交互に取ることが分かる.
                \[
                  a_n=
                  \begin{cases}
                    2 &(n\text{が奇数}),\\
                    -\dfrac12 &(n\text{が偶数}).
                  \end{cases}
                \]
                これも一般項の表し方である.逆数を用いた式は,この二つの場合を$(-1)^{n-1}$によってまとめたものである.
          % C2
          \item まず次の項だけを左辺に残すと,
                \[
                  a_{n+1}=-\frac12a_n+\frac92
                \]
                である.不動点は,
                \[
                  \alpha=-\frac12\alpha+\frac92
                  \quad\Longleftrightarrow\quad \alpha=3
                \]
                となる.\,$b_n=a_n-3$とおけば,
                \[
                  b_{n+1}=-\frac12b_n,\qquad b_1=1-3=-2
                \]
                なので,
                \[
                  b_n=-2\left(-\frac12\right)^{n-1},\qquad
                  a_n=3-2\left(-\frac12\right)^{n-1}
                \]
                を得る.負の公比により,\,$a_n-3$の符号が1項ごとに入れ替わる.
                例えば,\,$a_1=1<3$であり,\,$a_2=4>3$である.
          % U2
          \item $a_n+t b_n$が等比数列になるように$t$を探す.
                            \[
                              a_{n+1}+t b_{n+1}
                                =(4+t)a_n+(t-2)b_n.
                            \]
                                これが$(4+t)(a_n+t b_n)$と一致するためには,
                            \[
                              t-2=t(4+t),\qquad
                              (t+1)(t+2)=0
                            \]
                                であればよい.
                                $t=-1,-2$をそれぞれ用いると,
                            \[
                              a_{n+1}-b_{n+1}=3(a_n-b_n),
                            \]
                            \[
                              a_{n+1}-2b_{n+1}=2(a_n-2b_n).
                            \]
                                初項から,
                            \[
                              a_n-b_n=2\cdot3^{n-1},\qquad
                              a_n-2b_n=2^{n-1}.
                            \]
                                この2式の差を取ると$b_n=2\cdot3^{n-1}-2^{n-1}$である.
                                さらに$a_n=b_n+2\cdot3^{n-1}$より,
                            \[
                              a_n=4\cdot3^{n-1}-2^{n-1},\qquad
                              b_n=2\cdot3^{n-1}-2^{n-1}.
                            \]

                \noindent \textbf{【別解】}\\
                第1式から$2b_n=4a_n-a_{n+1}$である.
                                第2式を使うと,
                            \begin{align*}
                              a_{n+2}
                                &=4a_{n+1}-2b_{n+1}\\
                                &=4a_{n+1}-2(a_n+b_n)\\
                                &=5a_{n+1}-6a_n.
                            \end{align*}
                                $a_1=3,\ a_2=4\cdot3-2\cdot1=10$であり, 特性方程式は
                            \[
                              r^2-5r+6=(r-2)(r-3)=0.
                            \]
                                $a_n=C\cdot2^{n-1}+D\cdot3^{n-1}$とおくと,
                            \[
                              C+D=3,\qquad 2C+3D=10
                            \]
                                より$C=-1,\ D=4$である.
                                $\displaystyle b_n=\frac{4a_n-a_{n+1}}{2}$から$b_n$も求められる.
          % D2
          \item $a_n$を右辺にまとめると,
                \begin{align*}
                  2a_{n+1}&=2a_n+6,\\
                  a_{n+1}&=a_n+3
                \end{align*}
                となる.したがって,初項$-4$,\,公差$3$の等差数列であるから,
                \[
                  a_n=-4+3(n-1)=3n-7
                \]
                である.係数が両辺に分かれていても,まず隣り合う項の差を取り出せばよい.
          % G2
          \item 両辺から$a_{n+1}$を引き,さらに$5a_n$を加えると,
                \[
                  a_{n+1}=6a_n
                \]
                となる.定数項はなく,毎回同じ倍率で増える数列である.
                よって,初項$3$,\,公比$6$の等比数列として,
                \[
                  a_n=3\cdot6^{n-1}
                \]
                を得る.
          % R2
          \item $n\ge1$より$(n+1)(n+2)\ne0$であるから,
                \[
                  a_{n+1}=\frac{n}{n+1}\frac{n+3}{n+2}a_n
                \]
                と書ける.この2つの分数の積を分けて掛ければ,\,$n\ge2$のとき,
                \begin{align*}
                  a_n
                    &=6\prod_{k=1}^{n-1}\frac{k}{k+1}
                        \prod_{k=1}^{n-1}\frac{k+3}{k+2}\\
                    &=6\left(\frac12\frac23\cdots\frac{n-1}{n}\right)
                        \left(\frac43\frac54\cdots\frac{n+2}{n+1}\right)\\
                    &=6\cdot\frac1n\cdot\frac{n+2}{3}\\
                    &=\frac{2(n+2)}n
                \end{align*}
                となる.分子と次の分母が相殺するので,大きな積を計算する必要はない.
                $n=1$でも右辺は$6$であり,初項と一致する.

                \noindent \textbf{【別解】}\\
                両辺が同じ形になるように,
                \[
                  b_n=\frac{n}{n+2}a_n
                \]
                と置く.もとの漸化式より,
                \[
                  b_{n+1}
                    =\frac{n+1}{n+3}
                      \frac{n(n+3)}{(n+1)(n+2)}a_n
                    =\frac{n}{n+2}a_n
                    =b_n
                \]
                である.初項は$b_1=\dfrac13\cdot6=2$だから,\,$b_n=2$であり,
                \[
                  a_n=\frac{2(n+2)}n
                \]
                となる.
          % M2
          \item 二次式を引いて等比型にすることを考える.
                右辺は,
                \[
                  2n^2+2n+1=(n+1)^2+n^2
                \]
                と書ける.このため与式は,
                \[
                  a_{n+1}-(n+1)^2=-(a_n-n^2)
                \]
                となる.\,$b_n=a_n-n^2$とおけば,
                \[
                  b_{n+1}=-b_n,\qquad b_1=3-1=2
                \]
                なので,
                \[
                  b_n=2(-1)^{n-1},\qquad a_n=n^2+2(-1)^{n-1}
                \]
                を得る.二次式の部分$n^2$を引いた残りは,\,$2,-2,2,-2,\ldots$と交互に変わる.
          % S2
          \item $n=1$のとき与式は$S_1=a_1$となるだけなので, この式だけでは$a_1$は決まらない.
                                ただし, すべての正整数$n$で与式が成り立つことを使えば, 初項も決められる.
                            \begin{align*}
                              a_{n+1}
                                &=S_{n+1}-S_n\\
                                &=a_{n+1}+n2^n\\
                                &\quad{}-a_n-(n-1)2^{n-1}.
                            \end{align*}
                                両辺の$a_{n+1}$を消すと,
                            \begin{align*}
                              a_n
                                &=n2^n-(n-1)2^{n-1}\\
                                &=(n+1)2^{n-1}.
                            \end{align*}
                                この計算は$n\ge1$で成り立つので,\,$a_1=2$も含めて求められた.
          % L2
          \item 項の加減を含む式のまま対数を取っても簡単にならないので,まず平方完成する.
                \[
                  a_{n+1}-2=(a_n-2)^2
                \]
                ここで$a_1-2=3>0$であり,正の数の2乗は正なので,
                数学的帰納法により全ての$n$で$a_n-2>0$となる.
                そこで,\,$b_n=\log_3(a_n-2)$とおくと,
                \[
                  b_{n+1}=2b_n,\qquad b_1=\log_3 3=1
                \]
                である.従って,
                \[
                  b_n=2^{n-1},\qquad a_n-2=3^{\,2^{n-1}}
                \]
                なので,
                \[
                  a_n=2+3^{\,2^{n-1}}
                \]
                を得る.\,$n=1$でも$a_1=5$となる.

                \noindent \textbf{【別解：繰り返し2乗する】}\\
                \,$c_n=a_n-2$とおけば,\,$c_{n+1}=c_n^2,c_1=3$である.
                従って指数が$1,2,4,8,\ldots$と変わり,
                \[
                  c_n=3^{2^{n-1}}
                \]
                が予想できる.この式は$n=1$で成り立ち,
                \[
                  (3^{2^{n-1}})^2=3^{2^n}
                \]
                より次の項でも成り立つので,数学的帰納法で確認できる.
          % T3
          \item 特性方程式は
                \[
                  x^2-4x+4=(x-2)^2=0
                \]
                であり,重解$2$を持つ.二つの異なる補助数列は作れないので,
                \[
                  b_n=a_{n+1}-2a_n
                \]
                だけをまず調べる.
                \[
                  b_{n+1}=a_{n+2}-2a_{n+1}
                  =2(a_{n+1}-2a_n)=2b_n.
                \]
                $b_1=0-2\cdot1=-2$だから,
                \[
                  b_n=-2^n,\qquad a_{n+1}-2a_n=-2^n.
                \]
                両辺を$2^n$で割ると,
                \[
                  \frac{a_{n+1}}{2^n}-\frac{a_n}{2^{n-1}}=-1.
                \]
                $\displaystyle c_n=\frac{a_n}{2^{n-1}}$と置けば,\,$c_1=1,\,c_{n+1}=c_n-1$となる.よって,
                \[
                  c_n=1-(n-1)=2-n,
                  \qquad a_n=(2-n)2^{n-1}.
                \]
                割ったものは$2^n$であり,\,$a_2=0$でもこの操作に問題はない.

                \noindent \textbf{【別解】}\\
                初めから$\displaystyle c_n=\frac{a_n}{2^{n-1}}$と置き,もとの漸化式を$2^{n+1}$で割ると,
                \[
                  c_{n+2}-2c_{n+1}+c_n=0.
                \]
                したがって,
                \[
                  c_{n+2}-c_{n+1}=c_{n+1}-c_n.
                \]
                つまり,\,$\{c_n\}$は等差数列である.\,$c_1=1,\,c_2=0$より公差は$-1$なので,
                \[
                  c_n=2-n,\qquad a_n=(2-n)2^{n-1}.
                \]
                重解に対応する指数の部分を除くと,等差数列が残ることを直接確かめられる.
          % M3
          \item 加わる項が$-2^n$なので,同じ漸化式を満たす$g(n)=c\,2^n$を探す.
                \[
                  c\,2^{n+1}=3c\,2^n-2^n
                \]
                両辺を$2^n$で割ると$2c=3c-1$であり,\,$c=1$である.
                従って,\,$b_n=a_n-2^n$とおけば,
                \[
                  b_{n+1}=a_{n+1}-2^{n+1}=3(a_n-2^n)=3b_n
                \]
                となる.\,$b_1=1-2=-1$より,
                \[
                  b_n=-3^{n-1},\qquad a_n=2^n-3^{n-1}
                \]
                を得る.一般項は$n=1$で$1$となる.

                \noindent \textbf{【別解：指数列で割る】}\\
                \,$\displaystyle c_n=\frac{a_n}{2^{n-1}}$とおくと,
                \[
                  c_{n+1}=\frac32c_n-1,\qquad c_1=1
                \]
                である.不動点は$2$なので,
                \[
                  c_n-2=-\left(\frac32\right)^{n-1}
                \]
                となる.従って,
                \[
                  a_n=2^{n-1}\left\{2-\left(\frac32\right)^{n-1}\right\}
                     =2^n-3^{n-1}
                \]
                を得る.
          % U3
          \item 第2式はそれだけで解くことができる.
                                $b_1=2,\ b_{n+1}=2b_n$より,
                            \[
                              b_n=2^n.
                            \]
                                これを第1式に代入すると,
                            \[
                              a_{n+1}=2a_n+2^n.
                            \]
                                両辺を$2^n$で割り,
                            \[
                              \frac{a_{n+1}}{2^n}
                                =\frac{a_n}{2^{n-1}}+1
                            \]
                                と書き直す.
                                $\displaystyle c_n=\frac{a_n}{2^{n-1}}$とおけば$c_{n+1}=c_n+1,\ c_1=1$であるから,\,$c_n=n$である.
                                したがって,
                            \[
                              a_n=n2^{n-1},\qquad b_n=2^n.
                            \]

                \noindent \textbf{【別解】}\\
                $b_n=a_{n+1}-2a_n$を使い,\,$b_{n+1}=2b_n$から$b_n$を消去すると,
                            \[
                              a_{n+2}-2a_{n+1}
                                =2(a_{n+1}-2a_n).
                            \]
                                よって$a_{n+2}=4a_{n+1}-4a_n$である.
                                特性方程式は$(r-2)^2=0$である.
                                この形で$\displaystyle c_n=\frac{a_n}{2^{n-1}}$とおくと,
                            \[
                              c_{n+2}-2c_{n+1}+c_n=0
                            \]
                                となる. つまり$c_{n+1}-c_n$は一定なので,\,$c_n=C+D(n-1)$と書ける.
                                したがって,
                            \[
                              a_n=\{C+D(n-1)\}2^{n-1}.
                            \]
                                $a_1=1,\ a_2=2a_1+b_1=4$より$C=D=1$であり,\,$a_n=n2^{n-1}$を得る.
          % Q3
          \item 不動点の方程式は
                \[
                  x=\frac{3x+2}{x+2}
                  \quad\Longleftrightarrow\quad
                  x^2-x-2=(x-2)(x+1)=0
                \]
                なので,不動点は$2$と$-1$である.
                また,\,$a_1>2$であり,
                \[
                  a_{n+1}-2=\frac{a_n-2}{a_n+2}
                \]
                から,\,$a_n>2$なら$a_{n+1}>2$である.
                よって,すべての項で$a_n>2$となり,\,$a_n+2$と$a_n+1$は$0$ではない.
                そこで,
                \[
                  b_n=\frac{a_n-2}{a_n+1}
                \]
                と置くと,
                \begin{align*}
                  b_{n+1}
                  &=\cfrac{\cfrac{3a_n+2}{a_n+2}-2}
                           {\cfrac{3a_n+2}{a_n+2}+1}\\
                  &=\frac{a_n-2}{4a_n+4}
                  =\frac14b_n.
                \end{align*}
                $\displaystyle b_1=\frac{2}{5}$より,
                \[
                  b_n=\frac25\left(\frac14\right)^{n-1}.
                \]
                $b_n(a_n+1)=a_n-2$を$a_n$について解いて,
                \[
                  a_n=\frac{2+b_n}{1-b_n}
                  =\frac{10\cdot4^{n-1}+2}{5\cdot4^{n-1}-2}.
                \]
                $\displaystyle 0<b_n\leq\frac{2}{5}$なので,最後の変形で割った$1-b_n$も$0$ではない.

                \noindent \textbf{【別解】}\\
                不動点$2$からの差について,
                \[
                  a_{n+1}-2=\frac{a_n-2}{a_n+2}
                \]
                が成り立つ.\,$a_n>2$なので,この差の逆数を取ることができる.
                $\displaystyle c_n=\frac{1}{a_n-2}$と置くと,
                \[
                  c_{n+1}=\frac{a_n+2}{a_n-2}
                  =1+4c_n,\qquad c_1=\frac12.
                \]
                よって,
                \[
                  c_{n+1}+\frac13=4\left(c_n+\frac13\right),
                  \qquad
                  c_n=\frac56\cdot4^{n-1}-\frac13.
                \]
                したがって,
                \[
                  a_n=2+\frac1{c_n}
                  =2+\frac6{5\cdot4^{n-1}-2}
                  =\frac{10\cdot4^{n-1}+2}{5\cdot4^{n-1}-2}.
                \]
                二つの不動点からの比を使う方法と,一つの不動点からの差の逆数を使う方法は,同じ漸化式を別の形で単純にしている.
          % G3
          \item $n\ge1$より$n+1\ne0$である.両辺を$n+1$で割れば,
                \[
                  a_{n+1}=\frac13a_n
                \]
                となる.係数に$n$が含まれていても,整理した後の倍率は一定である.
                したがって,
                \[
                  a_n=4\left(\frac13\right)^{n-1}
                \]
                である.
          % D3
          \item 両辺から$n$を引けば,
                \[
                  a_{n+1}=a_n+\sqrt2
                \]
                となる.もとの式には$n$が現れるが,増加する量は$n$によらず一定である.
                よって,公差$\sqrt2$の等差数列として,
                \begin{align*}
                  a_n
                    &=1-\sqrt2+(n-1)\sqrt2\\
                    &=1+(n-2)\sqrt2
                \end{align*}
                を得る.公差が無理数でも,等差数列の一般項の求め方は同じである.
          % L3
          \item 全項が正であるので,\,$b_n=\log_2a_n$とおくことができる.
                すると,
                \[
                  b_{n+1}=\frac{n+2}{n}b_n,\qquad b_1=1
                \]
                となる.これは対数を取った後の数列が階比型となる例である.
                \,$n\ge2$では,
                \[
                  b_n=\prod_{k=1}^{n-1}\frac{k+2}{k}
                     =\frac{3\cdot4\cdots(n+1)}{1\cdot2\cdots(n-1)}
                     =\frac{n(n+1)}2
                \]
                である.最終式は$n=1$でも$b_1=1$を与える.
                従って,
                \[
                  a_n=2^{b_n}=2^{\frac{n(n+1)}2}
                \]
                となる.分数の指数が含まれていても,対数を取れば一次の係数になる.

                \noindent \textbf{【別解：対数を取らず指数をそろえる】}\\
                階比型の解法で現れた因子$n(n+1)$を利用し,
                \[
                  c_n=a_n^{\frac{2}{n(n+1)}}
                \]
                とおく.全項が正なので指数法則を使うことができ,
                \[
                  c_{n+1}
                  =\left(a_n^{\frac{n+2}{n}}\right)^{\frac{2}{(n+1)(n+2)}}
                  =a_n^{\frac{2}{n(n+1)}}=c_n
                \]
                となる.\,$c_1=2$だから$c_n=2$であり,
                \[
                  a_n=2^{\frac{n(n+1)}2}
                \]
                を得る.
          % S3
          \item 総和$S_n$も1つの数列なので, まずこれについて解く.
                            \[
                              S_{n+1}+1=3(S_n+1),\qquad S_1+1=2
                            \]
                                より,
                            \[
                              S_n=2\cdot3^{n-1}-1.
                            \]
                                $n\ge2$に対しては,
                            \begin{align*}
                              a_n
                                &=S_n-S_{n-1}\\
                                &=(2\cdot3^{n-1}-1)
                                  -(2\cdot3^{n-2}-1)\\
                                &=4\cdot3^{n-2}.
                            \end{align*}
                                一方, 初項は$a_1=S_1=1$である.
                                上で求めた式に$n=1$を代入すると$\displaystyle \frac{4}{3}$となり, 初項とは一致しない.
                                したがって,
                            \[
                              a_1=1,\qquad a_n=4\cdot3^{n-2}\quad(n\ge2).
                            \]

                \noindent \textbf{【別解】}\\
                $n\ge2$において,
                            \begin{align*}
                              a_{n+1}
                                &=S_{n+1}-S_n\\
                                &=(3S_n+2)-(3S_{n-1}+2)=3a_n.
                            \end{align*}
                                ただし,\,$S_n=3S_{n-1}+2$を使ったので, この式の導出は$n\ge2$で行っている.
                                $S_2=3S_1+2=5$より$a_2=5-1=4$であるから,
                            \[
                              a_n=4\cdot3^{n-2}\quad(n\ge2).
                            \]
                                $a_1=1$は別に付け加える.
          % F3
          \item $n\ge2$のとき,階差を順に加えると,
                \begin{align*}
                  a_n
                    &=3+2\sum_{k=1}^{n-1}3^k\\
                    &=3+2\frac{3(3^{n-1}-1)}{3-1}\\
                    &=3^n
                \end{align*}
                となる.最後の式は$n=1$でも初項$3$に一致する.
                階差が指数関数であっても,まず差を加え,その後で和を計算すればよい.

                \noindent \textbf{【別解1】}\\
                $2\cdot3^n=3^{n+1}-3^n$に気づけば,漸化式は
                \[
                  a_{n+1}-3^{n+1}=a_n-3^n
                \]
                と書ける.したがって,\,$a_n-3^n$は初項での値$a_1-3=0$に等しい.
                よって,\,$a_n=3^n$である.
                増加分が隣り合う指数関数の差になっているため,相殺を使っても解ける.

                \noindent \textbf{【別解2】}\\
                $a_n=3^n$を予想して帰納法を使うこともできる.
                初項は予想と一致し,\,$a_n=3^n$とすると,
                \[
                  a_{n+1}=3^n+2\cdot3^n=3^{n+1}
                \]
                となる.したがって,予想はすべての正整数$n$について成り立つ.
          % R3
          \item もとの式の両辺は,添字に合わせてその平方を掛けた同じ形になっている.
                そこで
                \[
                  b_n=n^2a_n
                \]
                と置くと,\,$b_{n+1}=b_n$である.
                初項は$b_1=1^2\cdot4=4$だから,すべての項で$b_n=4$となる.
                よって,
                \[
                  n^2a_n=4,\qquad a_n=\frac4{n^2}
                \]
                である.\,$n$は正整数なので,\,$n^2$で割ることができる.

                \noindent \textbf{【別解】}\\
                最初の数項は$4,\,1,\,\dfrac49,\,\dfrac14,\ldots$なので,\,$a_n=\dfrac4{n^2}$が予想できる.
                初項はこの式と一致する.
                $a_n=\dfrac4{n^2}$を仮定すると,
                \[
                  a_{n+1}
                    =\frac{n^2}{(n+1)^2}\frac4{n^2}
                    =\frac4{(n+1)^2}
                \]
                となる.したがって,数学的帰納法により一般項が確認できる.
          % C3
          \item 左辺が階差であっても,右辺には$a_n$が含まれる.
                そのまま既知の式の和を取ることはできないので,まず式を整理する.
                \[
                  a_{n+1}=\frac23a_n+2
                \]
                この式を満たす定数は$6$であり,
                \[
                  a_{n+1}-6=\frac23(a_n-6)
                \]
                となる.\,$b_n=a_n-6$とおくと,\,$b_1=-6$であるから,
                \[
                  b_n=-6\left(\frac23\right)^{n-1}
                \]
                よって,
                \[
                  a_n=6-6\left(\frac23\right)^{n-1}
                \]
                である.\,$n=1$では$0$となる.

                \noindent \textbf{【別解：増加量を求める】}\\
                与式を1つずらした式からもとの式を引くと,
                \[
                  (a_{n+2}-a_{n+1})-(a_{n+1}-a_n)
                  =-\frac13(a_{n+1}-a_n)
                \]
                となる.従って,\,$d_n=a_{n+1}-a_n$は,
                \[
                  d_{n+1}=\frac23d_n,\qquad d_1=2
                \]
                を満たす.これより,\,$n\ge2$では,
                \[
                  a_n=\sum_{k=1}^{n-1}2\left(\frac23\right)^{k-1}
                     =6\left\{1-\left(\frac23\right)^{n-1}\right\}
                \]
                を得る.この式は$n=1$でも正しい.
          % S4
          \item $S_1=a_1=1$であり,
                            \[
                              S_{n+1}=S_n+a_{n+1}=3S_n+2^n.
                            \]
                                そこで$c_n=S_n+2^n$とおくと,
                            \begin{align*}
                              c_{n+1}
                                &=3S_n+2^n+2^{n+1}\\
                                &=3(S_n+2^n)=3c_n.
                            \end{align*}
                                $c_1=1+2=3$より$c_n=3^n$となるので,
                            \[
                              S_n=3^n-2^n.
                            \]
                                よって$n\ge2$では,
                            \begin{align*}
                              a_n
                                &=S_n-S_{n-1}\\
                                &=2\cdot3^{n-1}-2^{n-1}.
                            \end{align*}
                                $n=1$でも右辺は$1=a_1$であるから,
                            \[
                              a_n=2\cdot3^{n-1}-2^{n-1}\quad(n\ge1).
                            \]

                \noindent \textbf{【別解】}\\
                $n\ge2$において, 与式から1つ前の式を引くと,
                            \[
                              a_{n+1}-a_n
                                =2(S_n-S_{n-1})+2^n-2^{n-1}
                                =2a_n+2^{n-1}.
                            \]
                                したがって$a_{n+1}=3a_n+2^{n-1}$である.
                                $a_2=2S_1+2=4=3a_1+1$より, この関係は$n=1$でも成り立つ.
                                $b_n=a_n+2^{n-1}$とおけば,
                            \[
                              b_{n+1}=3b_n,\qquad b_1=2.
                            \]
                                よって$b_n=2\cdot3^{n-1}$となり, 同じ答えを得る.
          % Q4
          \item $a_1>0$であり,\,$a_n>0$なら$\displaystyle a_{n+1}=2+\frac{3}{a_n}>0$なので,すべての項は正である.
                したがって,もとの分母$a_n$は$0$ではない.
                右辺を一つの分数にまとめると,
                \[
                  a_{n+1}=\frac{2a_n+3}{a_n}.
                \]
                不動点の方程式は
                \[
                  x=\frac{2x+3}{x}
                  \quad\Longleftrightarrow\quad
                  (x-3)(x+1)=0
                \]
                なので,\,$3$と$-1$が不動点である.
                $a_n+1>0$に注意して,
                \[
                  b_n=\frac{a_n-3}{a_n+1}
                \]
                と置けば,
                \begin{align*}
                  b_{n+1}
                  &=\frac{2a_n+3-3a_n}{2a_n+3+a_n}\\
                  &=-\frac{a_n-3}{3(a_n+1)}
                  =-\frac13b_n.
                \end{align*}
                $b_1=-1$だから,
                \[
                  b_n=-\left(-\frac13\right)^{n-1}
                  =\frac{(-1)^n}{3^{n-1}}.
                \]
                これを$a_n$に戻すと,
                \[
                  a_n=\frac{3+b_n}{1-b_n}
                  =\frac{3^n+(-1)^n}{3^{n-1}-(-1)^n}.
                \]
                最後の分母は,\,$n$が奇数なら$3^{n-1}+1$,偶数なら$3^{n-1}-1$である.
                偶数の場合は$n\geq2$なので,いずれも正であり$0$ではない.

                \noindent \textbf{【別解】}\\
                数列$\{u_n\}$を
                \[
                  u_1=u_2=1,\qquad u_{n+2}=2u_{n+1}+3u_n
                \]
                で定める.初めの二項が正で,右辺の係数も正なので,すべての$u_n$は正である.
                このとき,
                \[
                  \frac{u_{n+2}}{u_{n+1}}
                  =2+\cfrac{3}{\frac{u_{n+1}}{u_n}},
                  \qquad \frac{u_2}{u_1}=1.
                \]
                したがって,\,$\displaystyle a_n=\frac{u_{n+1}}{u_n}$となる.
                隣接3項間漸化式の特性方程式は
                \[
                  x^2-2x-3=(x-3)(x+1)=0
                \]
                である.\,$u_{n+2}+u_{n+1}=3(u_{n+1}+u_n)$に対応する等比数列と,
                $u_{n+2}-3u_{n+1}=-(u_{n+1}-3u_n)$に対応する等比数列を初項から求めると,
                \[
                  u_{n+1}+u_n=2\cdot3^{n-1},\qquad
                  u_{n+1}-3u_n=-2(-1)^{n-1}
                \]
                を得る.二式の差から,
                \[
                  u_n=\frac{3^{n-1}+(-1)^{n-1}}2.
                \]
                よって,
                \[
                  a_n=\frac{u_{n+1}}{u_n}
                  =\frac{3^n+(-1)^n}{3^{n-1}-(-1)^n}.
                \]
                項そのものではなく,隣り合う項の比として考えると,分数型と隣接3項間型がつながる.
          % G4
          \item 第$2$項から第$5$項までには$r$を$3$回掛けるから,
                \[
                  a_5=r^3a_2
                \]
                である.よって,\,$48=-6r^3$より$r^3=-8$となる.
                $r$は実数なので,\,$r=-2$である.
                さらに,\,$a_2=ra_1$より,
                \[
                  a_1=\frac{-6}{-2}=3
                \]
                となる.したがって,
                \[
                  a_n=3(-2)^{n-1}
                \]
                である.公比が負なので,正負が交互に現れることにも注意する.
          % F4
          \item 分母の2つの因数が$2$だけ離れていることに注目すると,
                \[
                  \frac{2}{(k+1)(k+3)}
                    =\frac1{k+1}-\frac1{k+3}
                \]
                と部分分数分解できる.よって,\,$n\ge2$のとき,
                \[
                  a_n=\sum_{k=1}^{n-1}
                       \left(\frac1{k+1}-\frac1{k+3}\right)
                \]
                である.和の範囲を分けて書くと,
                \begin{align*}
                  a_n
                    &=\sum_{j=2}^{n}\frac1j-\sum_{j=4}^{n+2}\frac1j\\
                    &=\frac12+\frac13-\frac1{n+1}-\frac1{n+2}\\
                    &=\frac56-\frac1{n+1}-\frac1{n+2}
                \end{align*}
                となる.2つの和は添字が$2$だけずれているので,両端にそれぞれ2項が残る.
                $n=2$では直接計算しても$a_2=\dfrac14$となり,得られた式と一致する.
                また,\,$n=1$を代入すると$0$となるから,この一般項は初項も含めて成り立つ.
          % T4
          \item 定数項を消すための定数解は,
                \[
                  \alpha=2\alpha+3\alpha-8
                  \quad\Longleftrightarrow\quad \alpha=2
                \]
                である.\,$b_n=a_n-2$と置けば,
                \[
                  b_{n+2}=2b_{n+1}+3b_n,
                  \qquad b_1=1,\quad b_2=3.
                \]
                特性方程式は
                \[
                  x^2-2x-3=(x-3)(x+1)=0.
                \]
                そこで,\,$c_n=b_{n+1}-3b_n$と置くと,
                \[
                  c_{n+1}=-c_n,\qquad c_1=3-3\cdot1=0.
                \]
                初項が$0$なので,\,$c_n=0$がすべての$n$で成り立つ.
                つまり,
                \[
                  b_{n+1}=3b_n.
                \]
                よって,\,$b_1=1$から$b_n=3^{n-1}$となり,
                \[
                  a_n=2+3^{n-1}.
                \]
                この問題では一方の補助数列が常に$0$になるため,二本を両方求める前に元の数列を等比数列に帰着できる.

                \noindent \textbf{【別解】}\\
                一般項を直接予想する前に,\,$a_{n+1}=3a_n-4$という短い漸化式が成り立つことを示す方法もある.
                まず$a_2=5=3a_1-4$である.
                ある$k$について$a_{k+1}=3a_k-4$が成り立つとき,
                \[
                  a_{k+2}-(3a_{k+1}-4)
                  =-a_{k+1}+3a_k-4
                  =0.
                \]
                したがって,数学的帰納法から$a_{n+1}=3a_n-4$が成り立つ.
                この一次の漸化式を
                \[
                  a_{n+1}-2=3(a_n-2)
                \]
                と変形すれば,\,$a_n=2+3^{n-1}$を得る.
          % U4
          \item $s_n=a_n+b_n,\ d_n=a_n-b_n$とおく.
                                2つの式を足すと定数項$1$と$-1$が消え,
                            \[
                              s_{n+1}=3s_n.
                            \]
                                一方, 引いたときには定数項の差は$2$となるので,
                            \[
                              d_{n+1}=d_n+2.
                            \]
                                $s_1=1,\ d_1=-1$より,
                            \[
                              s_n=3^{n-1},\qquad d_n=-1+2(n-1)=2n-3.
                            \]
                                したがって,
                            \begin{align*}
                              a_n&=\frac{s_n+d_n}{2}
                                   =\frac{3^{n-1}+2n-3}{2},\\
                              b_n&=\frac{s_n-d_n}{2}
                                   =\frac{3^{n-1}-2n+3}{2}.
                            \end{align*}
          % L4
          \item $a_1=1>0$であり,\,$a_n>0$なら$2\sqrt{a_n}>0$である.
                従って全ての項が正となる.
                \,$b_n=\log_2a_n$とおくと,
                \[
                  b_{n+1}=1+\frac12b_n,\qquad b_1=0
                \]
                である.不動点$2$との差を取れば,
                \[
                  b_{n+1}-2=\frac12(b_n-2)
                \]
                なので,
                \[
                  b_n=2-2\left(\frac12\right)^{n-1}=2-2^{2-n}
                \]
                となる.従って,
                \[
                  a_n=2^{\,2-2^{2-n}}
                \]
                を得る.\,$n=1$では指数が$0$となり,\,$a_1=1$である.

                \noindent \textbf{【別解：正の不動点で割る】}\\
                正の不動点は$4=2\sqrt4$を満たす$4$である.
                \,$\displaystyle c_n=\frac{a_n}{4}$とおくと,
                \[
                  c_{n+1}=\sqrt{c_n},\qquad c_1=\frac14
                \]
                となる.従って,
                \[
                  c_n=\left(\frac14\right)^{(\frac12)^{n-1}},\qquad
                  a_n=4\left(\frac14\right)^{(\frac12)^{n-1}}
                     =2^{\,2-2^{2-n}}
                \]
                を得る.1回進むたびに指数が半分になることを利用した.
          % C4
          \item 不動点を求めると,
                \[
                  \alpha=4\alpha-9\quad\Longleftrightarrow\quad\alpha=3
                \]
                である.\,$b_n=a_n-3$とおけば,
                \[
                  b_{n+1}=4b_n,\qquad b_1=t-3
                \]
                なので,
                \[
                  a_n=3+(t-3)4^{n-1}
                \]
                を得る.この計算では$t-3$で割っていないので,\,$t=3$の場合も含まれる.
                実際,\,$t=3$なら全ての$n$について$a_n=3$である.
                逆に定数数列なら$a_2=a_1$が必要であり,
                \[
                  4t-9=t\quad\Longleftrightarrow\quad t=3
                \]
                となる.従って求める値は$t=3$だけである.
          % R4
          \item 初項は正で,各段階の倍率$2^{2n+1}$も正なので,すべての項が正である.
                倍率を順に掛けると,\,$n\ge2$のとき,
                \begin{align*}
                  a_n
                    &=2\prod_{k=1}^{n-1}2^{2k+1}\\
                    &=2^{\,1+\sum_{k=1}^{n-1}(2k+1)}\\
                    &=2^{\,1+n(n-1)+(n-1)}\\
                    &=2^{n^2}
                \end{align*}
                となる.最後の式は$n=1$でも初項$2$に一致する.
                積の計算そのものではなく,指数の和を計算している点が大切である.

                \noindent \textbf{【別解1】}\\
                すべての項が正であるから,
                \[
                  b_n=\log_2 a_n
                \]
                と置くことができる.漸化式の両辺の対数を取ると,
                \[
                  b_{n+1}=b_n+2n+1,\qquad b_1=1
                \]
                となる.これは階差型であり,
                \[
                  b_n=1+\sum_{k=1}^{n-1}(2k+1)=n^2
                \]
                を得る.したがって,\,$a_n=2^{b_n}=2^{n^2}$である.
                指数の和を直接計算する方法と,対数を取る方法は,同じ計算を異なる形で表している.

                \noindent \textbf{【別解2】}\\
                $2n+1=(n+1)^2-n^2$を使うと,
                \[
                  \frac{a_{n+1}}{2^{(n+1)^2}}
                    =\frac{a_n}{2^{n^2}}
                \]
                となる.この比は初項で$\displaystyle \frac{2}{2}=1$だから,\,$a_n=2^{n^2}$を得る.
          % D4
          \item 第$3$項から第$8$項までには,公差$d$を$5$回加える.したがって,
                \[
                  a_8-a_3=5d
                \]
                であるから,\,$-7-8=5d$より$d=-3$を得る.
                また,\,$a_3=a_1+2d$であるから,
                \[
                  a_1=8-2(-3)=14
                \]
                である.よって,
                \[
                  a_n=14-3(n-1)=17-3n
                \]
                となる.求めた式に$n=3,\,8$を代入すれば,指定された2つの値をともに満たす.
          % M4
          \item $a_n$に掛かる係数と,加わる指数列の底がともに$2$である.
                そこで,\,$\displaystyle b_n=\frac{a_n}{2^{n-1}}$とおく.
                与式を$2^n$で割ると,
                \[
                  b_{n+1}=b_n+3,\qquad b_1=2
                \]
                となる.これは初項$2$,公差$3$の等差数列だから,
                \[
                  b_n=2+3(n-1)=3n-1
                \]
                である.よって,
                \[
                  a_n=(3n-1)2^{n-1}
                \]
                となる.変形後の式では不動点を求める必要がない.
                係数と底が等しいため,指数列で割った残りに$n$の一次式が現れる.
          % U5
          \item 第1式の添字を1つ進め, 第2式を代入すると,
                            \[
                              a_{n+2}=a_{n+1}+2b_{n+1}
                                =a_{n+1}+2a_n.
                            \]
                                $a_1=1,\ a_2=a_1+2b_1=1$である.
                                特性方程式は,
                            \[
                              r^2-r-2=(r-2)(r+1)=0.
                            \]
                                よって$a_n=C\cdot2^{n-1}+D(-1)^{n-1}$とおくと,
                            \[
                              C+D=1,\qquad 2C-D=1
                            \]
                                から$C=\frac{2}{3},\ D=\frac{1}{3}$となる.
                                したがって,
                            \[
                              a_n=\frac{2^n+(-1)^{n-1}}{3}.
                            \]
                                $n\ge2$では$b_n=a_{n-1}$であるから,
                            \[
                              b_n=\frac{2^{n-1}+(-1)^{n-2}}{3}
                                =\frac{2^{n-1}-(-1)^{n-1}}{3}.
                            \]
                                最後の式は$n=1$でも$0=b_1$となるので,
                            \[
                              b_n=\frac{2^{n-1}-(-1)^{n-1}}{3}\quad(n\ge1).
                            \]

                \noindent \textbf{【別解】}\\
                $c_n=a_n+b_n,\ d_n=a_n-2b_n$とおくと,
                            \[
                              c_{n+1}=2c_n,\qquad d_{n+1}=-d_n.
                            \]
                                $c_1=d_1=1$より$c_n=2^{n-1},\ d_n=(-1)^{n-1}$である.
                                $c_n-d_n=3b_n,\ 2c_n+d_n=3a_n$を用いると,\,$a_n,b_n$を直接求められる.
          % R5
          \item 係数$n$と$n+1$をそれぞれ階乗に吸収するため,
                \[
                  b_n=\frac{a_n}{(n-1)!\,n!}
                \]
                と置く.実際,
                \begin{align*}
                  b_{n+1}
                    &=\frac{a_{n+1}}{n!(n+1)!}\\
                    &=\frac{n(n+1)a_n}{2n!(n+1)!}\\
                    &=\frac12\frac{a_n}{(n-1)!\,n!}\\
                    &=\frac12b_n
                \end{align*}
                となる.ここでは$n!=n(n-1)!$,\, $(n+1)!=(n+1)n!$を用いた.
                $0!=1$より$b_1=1$なので,\,$b_n=2^{-(n-1)}$である.
                元に戻せば,
                \[
                  a_n=\frac{(n-1)!\,n!}{2^{n-1}}
                \]
                を得る.

                \noindent \textbf{【別解】}\\
                倍率を分けて掛ければ,
                \begin{align*}
                  a_n
                    &=\prod_{k=1}^{n-1}\frac{k(k+1)}2\\
                    &=\frac1{2^{n-1}}
                        \left(\prod_{k=1}^{n-1}k\right)
                        \left(\prod_{k=1}^{n-1}(k+1)\right)\\
                    &=\frac{(n-1)!\,n!}{2^{n-1}}
                \end{align*}
                となる.$n=1$では2つの積を空積$1$として扱うため,初項も同じ式に含まれる.
          % Q5
          \item 不動点の方程式は
                \[
                  x=2-\frac1x
                  \quad\Longleftrightarrow\quad
                  (x-1)^2=0
                \]
                であり,重解$1$を持つ.
                異なる二つの不動点を使った比は作れないので,\,$a_n-1$の変化を調べる.
                \[
                  a_{n+1}-1=1-\frac1{a_n}=\frac{a_n-1}{a_n}.
                \]
                $a_1>1$であり,\,$a_n>1$なら右辺は正なので$a_{n+1}>1$である.
                したがって,すべての項で$a_n>1$となり,もとの分母$a_n$も,差$a_n-1$も$0$ではない.
                両辺の逆数を取って,
                \[
                  \frac1{a_{n+1}-1}
                  =\frac{a_n}{a_n-1}
                  =1+\frac1{a_n-1}.
                \]
                $\displaystyle b_n=\frac{1}{a_n-1}$と置けば,\,$b_1=1,\,b_{n+1}=b_n+1$となる.
                よって,
                \[
                  b_n=n,\qquad a_n=1+\frac1n=\frac{n+1}{n}.
                \]
                不動点が重なる場合には,差の逆数から等差数列が現れた.

                \noindent \textbf{【別解】}\\
                数項を計算すると,
                \[
                  a_1=\frac21,\quad a_2=\frac32,\quad a_3=\frac43
                \]
                なので,\,$\displaystyle a_n=\frac{n+1}{n}$と予想できる.
                この式は$n=1$で成り立つ.\,$\displaystyle a_k=\frac{k+1}{k}$を仮定すると,
                \[
                  a_{k+1}=2-\frac{k}{k+1}
                  =\frac{k+2}{k+1}.
                \]
                したがって,数学的帰納法によって一般項を確かめられる.
          % G5
          \item $a_1=0$より,\,$a_2=7a_1=0$である.
                同様に,ある項が$0$なら次の項も$0$になるので,すべての項が$0$である.
                よって,
                \[
                  a_n=0
                \]
                を得る.等比型の一般項を直接用いても,\,$a_n=0\cdot7^{n-1}=0$となる.
                この例では$a_n=0$なので,\,$\displaystyle \frac{a_{n+1}}{a_n}$を作るために$a_n$で割ってはいけない.
          % L5
          \item 初項は正であり,正の項からは次の項も正となるので,全項について対数を取ることができる.
                \,$b_n=\log_2a_n$とおけば,
                \[
                  b_{n+1}=2b_n-n+1,\qquad b_1=2
                \]
                となる.項ごとに変わる$-n+1$を消すため,同じ式を満たす一次式$g(n)=un+v$を探す.
                \[
                  g(n+1)=2g(n)-n+1
                \]
                係数を比べると,
                \[
                  u=2u-1,\qquad u+v=2v+1
                \]
                なので,\,$u=1,v=0$であり,\,$g(n)=n$となる.
                従って,\,$c_n=b_n-n$とおけば,
                \[
                  c_{n+1}=2c_n,\qquad c_1=2-1=1
                \]
                である.よって,
                \[
                  b_n=n+2^{n-1},\qquad a_n=2^{\,n+2^{n-1}}
                \]
                を得る.指数に$n=1$を代入すると$2$なので,初項は$4$となる.

                \noindent \textbf{【別解：指数列で割ってから指数を追う】}\\
                対数を取った解法で$n$を引くことは,もとの項を$2^n$で割ることに対応する.
                \,$\displaystyle d_n=\frac{a_n}{2^n}$とおけば,
                \[
                  d_{n+1}=\frac{a_{n+1}}{2^{n+1}}
                         =\frac{a_n^2}{2^{2n}}=d_n^2,\qquad d_1=2
                \]
                となる.従って,
                \[
                  d_n=2^{2^{n-1}},\qquad
                  a_n=2^nd_n=2^{\,n+2^{n-1}}
                \]
                を得る.
          % F5
          \item $n\ge2$のとき,階差を加えると,
                \[
                  a_n=1+\sum_{k=1}^{n-1}(-1)^k
                \]
                となる.右辺の和は初項$-1$,\,公比$-1$の等比数列の和なので,
                \begin{align*}
                  a_n
                    &=1-\frac{1-(-1)^{n-1}}2\\
                    &=\frac{1+(-1)^{n-1}}2
                \end{align*}
                である.この式は$n=1$でも$1$となる.
                したがって,奇数番目の項は$1$,\,偶数番目の項は$0$である.

                \noindent \textbf{【別解】}\\
                2回分の変化をまとめると,
                \[
                  a_{n+2}
                    =a_n+(-1)^n+(-1)^{n+1}
                    =a_n
                \]
                となる.また,\,$a_1=1$,\, $a_2=0$であるから,数列は$1,\,0,\,1,\,0,\ldots$と繰り返す.
                この偶奇による違いを1つの式で表せば,
                \[
                  a_n=\frac{1+(-1)^{n-1}}2
                \]
                となる.
          % C5
          \item 最初の2回の更新より,
                \[
                  5=8p+q,\qquad \frac72=5p+q
                \]
                である.2式を引くと,
                \[
                  -\frac32=-3p
                \]
                なので,\,$p=\frac12$であり,\,$q=5-8\cdot\frac12=1$となる.
                従って,
                \[
                  a_{n+1}=\frac12a_n+1
                \]
                である.不動点$2$との差を取り,\,$b_n=a_n-2$とおくと,
                \[
                  b_{n+1}=\frac12b_n,\qquad b_1=6
                \]
                より,
                \[
                  a_n=2+6\left(\frac12\right)^{n-1}
                \]
                を得る.これに$n=1,2,3$を代入すると,順に$8,5,\frac72$となる.

                \noindent \textbf{【別解：階差から一般項を求める】}\\
                係数を決めた後,漸化式を1つずらして引くと,
                \[
                  d_n=a_{n+1}-a_n,\qquad
                  d_{n+1}=\frac12d_n,\qquad d_1=-3
                \]
                となる.従って,\,$n\ge2$では,
                \[
                  a_n=8-3\sum_{k=1}^{n-1}\left(\frac12\right)^{k-1}
                     =2+6\left(\frac12\right)^{n-1}
                \]
                を得る.最終式は$n=1$でも成り立つ.
          % T5
          \item すべての項を同じ定数$\alpha$に置き換えると,左辺は
                \[
                  \alpha-3\alpha+2\alpha=0
                \]
                となり,右辺の$1$とは一致しない.
                したがって,この問題には定数解がなく,定数を引くだけでは右辺を$0$にできない.
                そこで,階差
                \[
                  d_n=a_{n+1}-a_n
                \]
                を考える.漸化式を整理すると,
                \[
                  a_{n+2}-a_{n+1}=2(a_{n+1}-a_n)+1
                \]
                なので,
                \[
                  d_{n+1}=2d_n+1,\qquad d_1=2-1=1.
                \]
                これを$d_{n+1}+1=2(d_n+1)$と変形すれば,
                \[
                  d_n+1=2\cdot2^{n-1}=2^n,
                  \qquad d_n=2^n-1.
                \]
                よって,\,$n\geq2$では,
                \begin{align*}
                  a_n
                  &=a_1+\sum_{k=1}^{n-1}(a_{k+1}-a_k)\\
                  &=1+\sum_{k=1}^{n-1}(2^k-1)\\
                  &=1+(2^n-2)-(n-1)=2^n-n.
                \end{align*}
                この式は$n=1$でも$a_1=1$を与える.したがって,
                \[
                  a_n=2^n-n.
                \]
                階差もすぐに等比数列になるわけではないが,さらに$1$を加えることで等比数列になった.

                \noindent \textbf{【別解】}\\
                定数解の代わりに,\,$-n$が漸化式を満たすことに気付けば,
                \[
                  -(n+2)-3\{-(n+1)\}+2(-n)=1
                \]
                なので,\,$b_n=a_n+n$と置くことができる.
                \[
                  b_{n+2}-3b_{n+1}+2b_n=0,
                  \qquad b_1=2,\quad b_2=4.
                \]
                ここで,\,$c_n=b_{n+1}-2b_n$と置くと,
                \[
                  c_{n+1}=c_n,\qquad c_1=4-2\cdot2=0.
                \]
                したがって,\,$b_{n+1}=2b_n$となり,\,$b_n=2^n$を得る.
                よって,
                \[
                  a_n=b_n-n=2^n-n.
                \]
                右辺の定数を消すために,項ごとに変わる一次式を引くこともできる.
          % D5
          \item 両辺を$4$で割ると,\,$a_{n+1}-a_n=0$である.つまり,次の項は前の項と等しい.
                初項が$\dfrac74$なので,
                \[
                  a_1=a_2=a_3=\cdots=\frac74
                \]
                となり,\,$a_n=\dfrac74$である.
                これは公差$0$の等差数列であり,公式を使っても
                \[
                  a_n=\frac74+(n-1)\cdot0=\frac74
                \]
                と求まる.
          % M5
          \item $(-1)^n$は$0$にならないので,\,$\displaystyle b_n=\frac{a_n}{(-1)^{n-1}}$とおける.
                与式を$(-1)^n$で割ると,
                \[
                  b_{n+1}=-2b_n+3,\qquad b_1=0
                \]
                となる.この式の不動点は$1$であり,
                \[
                  b_{n+1}-1=-2(b_n-1)
                \]
                なので,
                \[
                  b_n=1-(-2)^{n-1}
                \]
                を得る.もとの数列に戻すと,
                \[
                  a_n=(-1)^{n-1}\{1-(-2)^{n-1}\}
                     =(-1)^{n-1}-2^{n-1}
                \]
                となる.最後の等号には$(-1)^{n-1}(-2)^{n-1}=2^{n-1}$を用いた.

                \noindent \textbf{【別解：交代する項を直接引く】}\\
                \,$g_n=(-1)^{n-1}$とおくと,
                \[
                  g_{n+1}=2g_n+3(-1)^n
                \]
                が成り立つ.与式からこれを引けば,
                \[
                  a_{n+1}-g_{n+1}=2(a_n-g_n)
                \]
                となる.\,$a_1-g_1=-1$だから,
                \[
                  a_n-g_n=-2^{n-1}
                \]
                であり,同じ一般項を得る.
          % S5
          \item 与式の$n$を$n+1$に置き換えた式との差を取ると,
                            \[
                              a_{n+1}
                                =\frac{n+3}{3}a_{n+1}
                                 -\frac{n+2}{3}a_n.
                            \]
                                整理して,
                            \[
                              na_{n+1}=(n+2)a_n,\qquad
                              a_{n+1}=\frac{n+2}{n}a_n.
                            \]
                                $a_1=2$よりすべての項は正であり,\,$n\ge2$では,
                            \begin{align*}
                              a_n
                                &=2\prod_{k=1}^{n-1}\frac{k+2}{k}\\
                                &=2\frac{3\cdot4\cdots(n+1)}{1\cdot2\cdots(n-1)}\\
                                &=n(n+1).
                            \end{align*}
                                この式は$n=1$でも$a_1=2$と一致する.
                                したがって,
                            \[
                              a_n=n(n+1).
                            \]

                \noindent \textbf{【別解】}\\
                $na_{n+1}=(n+2)a_n$を
                            \[
                              \frac{a_{n+1}}{(n+1)(n+2)}
                                =\frac{a_n}{n(n+1)}
                            \]
                                と書き直す. よって$\displaystyle \frac{a_n}{n(n+1)}$は一定であり, 初項からその値は$\displaystyle \frac{2}{1\cdot2}=1$である.
                                したがって$a_n=n(n+1)$となる.
        \end{enumerate}

